% ============================================================
% 正文：Introduction
% 作者：Kehui Chen, Jacky Keung, Xiaoxue Ma
% 主题：MCP 故障定位（fault localization）根因实证研究
% ============================================================

\section{Introduction}\label{sec:intro}

\noindent Large language models (LLMs) are increasingly deployed in enterprise software
engineering, yet they struggle with cross-system invocation and multi-tool collaboration
\cite{hou2025model}\cite{yang2025survey}. To tackle this, Anthropic released the Model
Context Protocol (MCP) in late 2024 \cite{anthropic_mcp_2024}. MCP acts as a middle-layer
protocol that connects LLMs to external resources through four core components: Host,
Client, Server, and Resource \cite{yang2025survey}. It standardises tool invocation and
context exchange across heterogeneous systems, and has become a foundation for many
LLM-driven workflows in industry.

Existing MCP research largely falls into three directions: benchmark construction for
real-world LLM evaluation \cite{liu2025mcpeval}\cite{yang2025mcpsecbench}, system
integration and performance optimisation \cite{xu2025towards}\cite{li2025energyplus}, and
security-threat analysis with countermeasures \cite{hasan2025model}\cite{mcpSecurity2025}.
Most benchmarks focus on adversarial security attacks or high-level functional evaluation,
while comparatively little attention has been paid to \emph{ordinary runtime faults} and
the \emph{observability issues} that arise in real-world MCP deployments.

Despite these technical advances, empirical evidence on how MCP behaves in production
remains scarce. In particular, practitioners consistently report that \emph{it is
extremely difficult to localise the root cause of a fault} in an MCP pipeline. The
difficulty is not merely a subjective complaint: it stems from the way MCP SDKs report
errors. A single call runs between the Host, Client, and Server components, and a fault may
arise in any of them or during the transport of the messages between them; yet the layer at
which the resulting error is \emph{displayed} often differs from the component in which the
fault actually \emph{originated}. Moreover, some erroneous inputs are silently accepted by
the SDK and return a successful status without any error flag, making these defects
invisible to ordinary monitoring pipelines.

Existing empirical evidence on MCP adoption originates mostly from practitioner
interviews \cite{ma2025practitioners} or open-source issue mining. While interviews capture
subjective industrial pain points, they cannot provide \emph{controlled, reproducible}
technical evidence that explains \emph{why} certain problems occur. To our best knowledge,
there is still no end-to-end, runnable fault dataset for MCP runtime failure analysis. This
creates a gap between perceived real-world difficulties and verifiable technical root
causes.

To narrow this gap, we adopt a \emph{sequential-exploratory mixed-methods} research design.
In the first phase, we obtain direct insights from enterprise practitioners through
semi-structured interviews to identify the core pain points of MCP deployment, with special
attention to fault diagnosis. In the second phase, guided by these pain points, we mine a
\textbf{real-bug dataset} --- a complete census of the official
\verb|modelcontextprotocol/python-sdk| issue tracker (from 1,343 pure issues to 398
confirmed real bugs to 342 active after removing the 56 already
fixed in the pinned SDK), reproduce \textbf{49} of them at runtime over the production
Streamable-HTTP transport, and run observability experiments that empirically examine the
technical roots of fault-localisation difficulty. This triangulation connects qualitative
practitioner insight with objective experimental observation, and lets us quantify why fault
localisation in MCP is so hard.

Concretely, this work addresses the following research questions:
\begin{itemize}
  \item \textbf{RQ1} What types of runtime faults can occur across the four-layer MCP
        fault taxonomy (Host, Client, Server, Transport)?
  \item \textbf{RQ2} To what extent does MCP SDK error reporting suffer from
        \emph{manifestation-root-cause mismatch}, i.e.\ errors surfaced at one layer while
        the true root cause lies in another?
  \item \textbf{RQ3} To what extent does MCP SDK error reporting exhibit \emph{silent
        failures}, i.e.\ erroneous inputs that return a successful status without an error
        flag?
  \item \textbf{RQ4} What are the technical root causes that make MCP fault localisation
        difficult, and how can automated diagnostic guidance alleviate them?
\end{itemize}

The main contributions of this paper are summarised as follows:
\begin{enumerate}
  \item We present interview evidence of MCP adoption from 20 real-world practitioners
        across Internet and FinTech enterprises, characterising fault-diagnosis as the
        leading operational obstacle (Section~\ref{sec:results}).
  \item We define a four-layer MCP runtime-fault taxonomy (Host scheduling, Client message,
        Server tool, Transport) and build a \textbf{real-bug dataset} by mining the entire
        issue tracker of the official \verb|modelcontextprotocol/python-sdk|, reproducing
        \textbf{49} genuine bugs at runtime over the production Streamable-HTTP transport
        (Sections~\ref{sec:method},~\ref{sec:setup}).
  \item We empirically derive a \emph{three-channel exposure model} characterising how MCP
        faults actually surface: \emph{pushed-to-Server} (7/13 non-Server bugs),
        \emph{silent failure} (11/49 = 22.4\%), and \emph{caller-exception} (3/13); a further
        17/49 = 34.7\% are \emph{genericised}, their real root cause replaced by a generic
        message. None of the 13 non-Server bugs is localised at its true origin (10/13 = 76.9\%
        generalised mismatch).
  \item We build a genuine LLM diagnosis baseline: a model (DeepSeek-V4 via a CityU OpenAI
        proxy) that reads only the true SDK error text and predicts the root-cause layer
        achieves 77.6\% overall but only \textbf{15.4\%} on the 13 non-Server bugs (0\% on both
        Host- and Client-origin faults, misattributed to the Server layer), and the non-Server
        Top-K is only 23.1\% at Top-1 (84.6\% at Top-3).
  \item We design a \emph{runnable diagnostic-clue tool} that classifies the exposure channel
        and emits directional guidance without revealing the root-cause layer, and show it
        raises the non-Server localisation accuracy from 15.4\% to 23.1\% (+7.7 pp), whereas a
        bare channel-label control does not help (0\%), evidencing the value of honest,
        non-leaking diagnostic guidance.
  \item We run a repair experiment and find the clue \emph{does not} change repair success
        (69.6\% with or without it on the auto-verifiable subset), because with the full source
        the LLM reads the fault directly --- corroborating that localisation, not code-fixing,
        is the bottleneck in MCP debugging.
  \item We release the fault dataset, source scripts, and raw execution records, and derive
        actionable observability recommendations for future MCP protocol and SDK design.
\end{enumerate}

The remainder of this paper is organised as follows.
Section~\ref{sec:related} reviews related work on LLM tool integration, MCP background, and
existing fault-taxonomy research.
Section~\ref{sec:method} describes our mixed-methods methodology and defines the four-layer
fault taxonomy.
Section~\ref{sec:setup} details the real-bug dataset construction, the four-layer taxonomy
and the controlled reproduction / observability experiment configuration.
Section~\ref{sec:results} reports both the demand-side interview evidence and the results
of the fault-injection and diagnosis experiments.
Section~\ref{sec:discussion} discusses research and practical implications, together with
threats to validity.
Finally, Section~\ref{sec:concl} concludes the paper.